\section{AIME II 2006}
        \begin{enumerate}
        	\item (\textit{AIME II 2006}) In convex hexagon $ABCDEF$, all six sides are congruent, $\angle A$ and $\angle D$ are right angles, and $\angle B, \angle C, \angle E,$ and $\angle F$ are congruent. The area of the hexagonal region is $2116(\sqrt{2}+1).$ Find $AB$.


 
 \textbf{Answer: }046



\textbf{Solution 1}

Let the side length be called $x$, so $x=AB=BC=CD=DE=EF=AF$. 

 
                    \begin{figure}[H]
                        \centering
                        \includegraphics[width=0.5\textwidth]{Solution/AIME_II_2006_1_Solution_Figure_1.png}
                    \end{figure}
                    

 The diagonal $BF=\sqrt{AB^2+AF^2}=\sqrt{x^2+x^2}=x\sqrt{2}$. Then the areas of the triangles AFB and CDE in total are $\frac{x^2}{2}\cdot 2$,and the area of the rectangle BCEF equals $x\cdot x\sqrt{2}=x^2\sqrt{2}$

 Then we have to solve the equation 

 
 \[2116(\sqrt{2}+1)=x^2\sqrt{2}+x^2\]
 \[2116(\sqrt{2}+1)=x^2(\sqrt{2}+1)\]
 \[2116=x^2\]
 \[x=46\]

 Therefore, $AB$ is $\boxed{046}$.

 





\textbf{Solution 2}

Because $\angle B$, $\angle C$, $\angle E$, and $\angle F$ are congruent, the degree-measure of each of them is ${{720-2\cdot90}\over4}= 135$. Lines $BF$ and $CE$ divide the hexagonal region into two right triangles and a rectangle. Let $AB=x$. Then $BF=x\sqrt2$. Thus
 \begin{align*} 2116(\sqrt2+1)&=[ABCDEF]\\ &=2\cdot {1\over2}x^2+x\cdot x\sqrt2=x^2(1+\sqrt2), \end{align*}
so $x^2=2116$, and $x=\boxed{046}$.

 
                    \begin{figure}[H]
                        \centering
                        \includegraphics[width=0.5\textwidth]{Solution/AIME_II_2006_1_Solution_Figure_2.png}
                    \end{figure}
                    

 





\textbf{Solution 3}

Extend lines $AB$ and $CD$ to intersect at $G$ and $AF$ and $DE$ to intersect at $H$. Angles $B$, $C$, $E$, and $F$ are 135 degrees, so minor angles $B$, $C$, $E$, $F$ apart of triangles $G$ and $H$ respectively are 45 degrees. Let $GB = GC = HE = HF = a$, so we have that $AB = BC = \dots = FA = a\sqrt{2}$. The two side lengths of quadrilateral $AGDH$ are equal and the measure in degrees of all figures are 90, so the quadrilateral $AGDH$ is a square. It follows that the side length of this square is $a + a\sqrt{2} = a(\sqrt{2}+1)$, and the area of this figure is $(a(\sqrt{2}+1))^2$. The area of the hexagon is then $(a(\sqrt{2}+1))^2 - a^2$, so we have $a^2((\sqrt{2}+1)^2 - 1) = 2116(\sqrt{2}+1)$, so we have $2a^2(\sqrt{2}+1) = 2116(\sqrt{2}+1)$, then $\sqrt{2}a = 46$, so $a = \frac{46}{\sqrt{2}}$. The side length is $a\sqrt{2}$, so the answer is $\frac{46}{\sqrt{2}} \cdot \sqrt{2} = \boxed{046}$.

 




	\item (\textit{AIME II 2006}) The lengths of the sides of a triangle with positive area are $\log_{10} 12$, $\log_{10} 75$, and $\log_{10} n$, where $n$ is a positive integer. Find the number of possible values for $n$.


 
 \textbf{Answer: }893



\textbf{Solution}

By the \href{https://artofproblemsolving.com/wiki/index.php?title=Triangle\_Inequality}{Triangle Inequality} and applying the well-known logarithmic property $\log_{c} a + \log_{c} b = \log_{c} ab$, we have that

 $\log_{10} 12 + \log_{10} n > \log_{10} 75$

 $\log_{10} 12n > \log_{10} 75$

 $12n > 75$ 

 $n > \frac{75}{12} = \frac{25}{4} = 6.25$

 Also, 

 $\log_{10} 12 + \log_{10} 75 > \log_{10} n$

 $\log_{10} 12\cdot75 > \log_{10} n$

 $n < 900$ 

 Combining these two inequalities: 

 
 \[6.25 < n < 900\]

 Thus $n$ is in the set $(6.25 , 900)$; the number of positive integer $n$ which satisfies this requirement is $\boxed{893}$.

 





\textbf{Solution 1 Clarified}

Triangle Inequality on the upper bound of $n$ yields $n < 900$, which is when $\log(n) > \log(75), \log(12)$, say case 1. For the lower bound, we have that $\log(12) < \log(75)$ so case 2 is when $\log(75)$ is the longest side, hence $n > 6.25$. Otherwise, we cannot have $n > \frac{12}{75}$ as the lower bound because $\log(12)$ can never be the longest side of a triangle with another side $\log(75)$. Hence we continue by counting integral values $n \in [7,899]$ which is $\boxed{893}$.

 




	\item (\textit{AIME II 2006}) Let $P$ be the product of the first 100 positive odd integers. Find the largest integer $k$ such that $P$ is divisible by $3^k$.


 
 \textbf{Answer: }049



\textbf{Solution}

Note that the product of the first $100$ positive odd integers can be written as $1\cdot 3\cdot 5\cdot 7\cdots 195\cdot 197\cdot 199=\frac{1\cdot 2\cdots200}{2\cdot4\cdots200} = \frac{200!}{2^{100}\cdot 100!}$

 Hence, we seek the number of threes in $200!$ decreased by the number of threes in $100!.$

 There are 

 $\left\lfloor \frac{200}{3}\right\rfloor+\left\lfloor\frac{200}{9}\right\rfloor+\left\lfloor \frac{200}{27}\right\rfloor+\left\lfloor\frac{200}{81}\right\rfloor =66+22+7+2=97$

 threes in $200!$ and 

 $\left\lfloor \frac{100}{3}\right\rfloor+\left\lfloor\frac{100}{9}\right\rfloor+\left\lfloor \frac{100}{27}\right\rfloor+\left\lfloor\frac{100}{81}\right\rfloor=33+11+3+1=48$

 threes in $100!$

 Therefore, we have a total of $97-48=\boxed{049}$ threes.

 For more information, see also \href{https://artofproblemsolving.com/wiki/index.php?title=Factorial\#Prime\_factorization}{ prime factorizations of a factorial}.

 


 





\textbf{Solution 2}

We count the multiples of $3^k$ below 200 and subtract the count of multiples of $2\cdot 3^k$:

 
 \[\left\lfloor \frac{200}{3}\right\rfloor - \left\lfloor \frac{200}{6}\right\rfloor +\left\lfloor\frac{200}{9}\right\rfloor - \left\lfloor \frac{200}{18}\right\rfloor +\left\lfloor \frac{200}{27}\right\rfloor - \left\lfloor \frac{200}{54}\right\rfloor+\left\lfloor\frac{200}{81}\right\rfloor - \left\lfloor \frac{200}{162}\right\rfloor\]
 \[= 66 - 33 + 22 - 11 + 7 - 3 + 2 - 1 = 49.\]

 





\textbf{Solution 3}

We can use a modified version of Legendre's Formula. First, we count the number of multiples of 3 in the sequence 1, 3, 5, 7, 9, ..., 195, 197, 199. 

 This is the same as the number of multiples of 3 in the sequence 3, 9, 15, 21, ..., 192, 195. There are clearly 33 terms in this sequence. 

 Next, we count the number of multiples of 9 in the sequence 1, 3, 5, 7, 9, ..., 195, 197, 199. This is the same as the number of multiples of 9 in the sequence 9, 18, 27, 36, ...., 189, 198 - but there's a catch. Note that every other member of this sequence isn't odd and thus is not part of the product of the first 100 odd integers, so our new sequence is actually 9, 27, 45...189. Divide every term by 9 to get a new sequence; 1, 3, 5...21, which clearly has 11 terms. 

 Next, we similarly count the number of multiples of 27 in the sequence 1, 3, 5, 7, 9, ..., 195, 197, 199. This is just 27, 81, 135, 189, so 4 multiples here. 

 Finally, we count the number of multiples of 81 in the sequence 1, 3, 5, 7, 9, ..., 195, 197, 199. There is only one such multiple, 81. 

 Any power of 3 above 81 doesn't fit into our sequence. 

 Finally, we have 33+11+4+1=49.

 Our final answer is 49.

 (Note that I oversimplified this a lot, in real life we wouldn't have to list out the sequences as tediously as I did).

 





\textbf{Solution 4}

We are obviously searching for multiples of three set S of odd numbers 1-199. Starting with 3, every number $\equiv 2 \pmod{3}$ in set S will be divisible by 3. In other words, every number $\equiv 3 \pmod{6}$. This is because the LCM must be divisible by 3, and 2, because the set is comprised of only odd numbers. Using some simple math, there are 33 numbers that fit this description. All you have to do is find the largest odd number that is divisible by 3 and below 200, then add 3 and divide by 6. Next, we find the number of odd digits below 200 that are divisible by 9. The same strategy works, and gives us 11. 27 gives 3, and 81 gives 1. $33 + 11 + 4 + 1 = \boxed{49}$.

 
-jackshi2006

 





\textbf{Solution 5 (simpler solution 4)}

Consider all integers from 1 - 200. Because $3 \ \not | \ 200$, we just count number of integers from 1 - 200 divisble by 3. Half of those would be even, so we divide by 2, and we have the odd numbers with at least one power of 3. Between all odd numbers 1 - 200, we have 33 factors of 3. From 1 - 200, we have 11 factors of 9. From 1 - 200, we have 4 powers of 27 (found as $200 \pmod{27} = 11 < 13 \pmod{27}$), and 1 power of 81. So, the count total is $33 + 11 + 4 + 1 = \boxed{049}$.

 




	\item (\textit{AIME II 2006}) Let $(a_1,a_2,a_3,\ldots,a_{12})$ be a permutation of $(1,2,3,\ldots,12)$ for which


 $a_1>a_2>a_3>a_4>a_5>a_6 \mathrm{\  and \ } a_6<a_7<a_8<a_9<a_{10}<a_{11}<a_{12}.$An example of such a permutation is $(6,5,4,3,2,1,7,8,9,10,11,12).$ Find the number of such permutations.


 
 \textbf{Answer: }462



\textbf{Solution}

Clearly, $a_6=1$.  Now, consider selecting $5$ of the remaining $11$ values.  Sort these values in descending order, and sort the other $6$ values in ascending order.  Now, let the $5$ selected values be $a_1$ through $a_5$, and let the remaining $6$ be $a_7$ through ${a_{12}}$.  It is now clear that there is a \href{https://artofproblemsolving.com/wiki/index.php?title=Bijection}{bijection} between the number of ways to select $5$ values from $11$ and ordered 12-tuples $(a_1,\ldots,a_{12})$.  Thus, there will be ${11 \choose 5}=\boxed{462}$ such ordered 12-tuples.

 





\textbf{Solution 2}

There are $\binom{12}{6}$ ways to choose 6 numbers from $(1,2,3,\ldots,12)$, and then there will only be one way to order them. And since that $a_6<a_7$, only half of the choices will work, so the answer is $\frac{\binom{12}{6}}{2}=462$ 12-tuples - mathleticguyyyI’maSTAY

 





\textbf{Solution 3}

Clearly, $a_6=1$, and either $a_1$ or $a_{12}$ is 12.

 
Case 1: $a_1 = 12$

 In this case, there are 4 empty spaces between $a_1$ and $a_6$, and 6 empty spaces between $a_6$ and $a_{12}$. $\binom{10}{4}$ is 210. This splits the remaining 10 numbers into two distinct sets that are automatically ordered. For this reason, there is no need to multiply by two to count doubles or treat as a permutation.

 
Case 2: $a_{12} = 12$

 In this case, there are 5 empty spaces between $a_1$ and $a_6$, and 5 empty spaces between $a_6$ and $a_{12}$. $\binom{10}{5}$ is 252. Like last time, this splits the remaining 10 numbers into two distinct sets that are automatically ordered. It is important to realize that the two sets are distinct because one side has 12 and the other does not. There is no need to multiply by two.

 $210 + 252 = \boxed{462}$ ordered 12-tuples.

 
-jackshi2006

 




	\item (\textit{AIME II 2006}) When rolling a certain unfair six-sided die with faces numbered $1, 2, 3, 4, 5$, and $6$, the probability of obtaining face $F$ is greater than $\frac{1}{6}$, the probability of obtaining the face opposite is less than $\frac{1}{6}$, the probability of obtaining any one of the other four faces is $\frac{1}{6}$, and the sum of the numbers on opposite faces is 7. When two such dice are rolled, the probability of obtaining a sum of 7 is $\frac{47}{288}$. Given that the probability of obtaining face $F$  is $\frac{m}{n},$ where $m$ and $n$ are relatively prime positive integers, find $m+n.$


 
 \textbf{Answer: }029



\textbf{Solution 1}

\href{https://artofproblemsolving.com/wiki/index.php?title=Without\_loss\_of\_generality}{Without loss of generality}, assume that face $F$ has a 6, so the opposite face has a 1.  Let $A(n)$ be the probability of rolling a number $n$ on one die and let $B(n)$ be the probability of rolling a number $n$ on the other die. 7 can be obtained by rolling a 2 and 5, 5 and 2, 3 and 4, or 4 and 3. Each has a probability of $\frac{1}{6} \cdot \frac{1}{6} = \frac{1}{36}$, totaling $4 \cdot \frac{1}{36} = \frac{1}{9}$. Subtracting all these probabilities from $\frac{47}{288}$ leaves $\frac{15}{288}=\frac{5}{96}$ chance of getting a 1 on die $A$ and a 6 on die $B$ or a 6 on die $A$ and a 1 on die $B$:

 
 \[A(6)\cdot B(1)+B(6)\cdot A(1)=\frac{5}{96}\]

 Since the two dice are identical, $B(1)=A(1)$ and $B(6)=A(6)$ so

 
 \begin{align*}A(1)\cdot A(6)+A(1)\cdot A(6)&=\frac{5}{96}\\ A(1)\cdot A(6)&=\frac{5}{192}\end{align*}


 Also, we know that $A(2)=A(3)=A(4)=A(5)=\frac{1}{6}$ and that the total probability must be $1$, so:

 
 \[A(1)+4 \cdot \frac{1}{6}+A(6)=\frac{6}{6} \Longrightarrow A(1)+A(6)=\frac{1}{3}\]

 Combining the equations:

 
 \begin{align*}A(6)\left(\frac{1}{3}-A(6)\right)&=\frac{5}{192}\\ 0 &= 192 \left(A(6)\right)^2 - 64 \left(A(6)\right) + 5\\ A(6)&=\frac{64\pm\sqrt{64^2 - 4 \cdot 5 \cdot 192}}{2\cdot192} =\frac{5}{24}, \frac{1}{8}\end{align*}
We know that $A(6)>\frac{1}{6}$, so it can't be $\frac{1}{8}$. Therefore, the probability is $\frac{5}{24}$ and the answer is $5+24=\boxed{29}$.

 Note also that the initial assumption that face $F$ was the face labelled 6 is unnecessary -- we would have carried out exactly the same steps and found exactly the same probability no matter which face it was. We could have labelled $A(6)$ as $p$, for example, and replaced the others with variables too, but the notation would have been harder to follow.

 





\textbf{Solution 2}

We have that the cube probabilities to land on its faces are $\frac{1}{6}$, $\frac{1}{6}$, $\frac{1}{6}$, $\frac{1}{6}$ ,$\frac{1}{6}+x$ ,$\frac{1}{6}-x$we also know that the sum could be 7 only when the faces in each of the two tosses are opposite hence the probability to get a 7 is:
 \[4 \cdot \left(\frac{1}{6} \right)^2+2 \left(\frac{1}{6}+x \right) \left(\frac{1}{6}-x \right)=\frac{47}{288}\]multiplying by 288 we get:
 \[32+16(1-6x)(6x+1)=47 \Longrightarrow 16(1-36x^2)=15\]dividing by 16 and rearranging we get:
 \[\frac{1}{16}=36x^2 \longrightarrow x=\frac{1}{24}\]so the probability F  which is greater than $\frac{1}{6}$ is equal $\frac{1}{6}+\frac{1}{24}=\frac{5}{24}\longrightarrow 24+5=\boxed{29}$

 





\textbf{Solution 3 (Alcumus)}

Let $p(a,b)$ denote the probability of obtaining $a$ on the first die and $b$ on the second. Then the probability of obtaining a sum of 7 is
 \[p(1,6)+p(2,5)+p(3,4)+p(4,3)+p(5,2)+p(6,1).\]Let the probability of obtaining face $F$ be $(1/6)+x$. Then the probability of obtaining the face opposite face $F$ is $(1/6)-x$. Therefore

 
 \begin{align*} {{47}\over{288}}&= 4\left({1\over6}\right)^2+2\left({1\over6}+x\right) \left({1\over6}-x\right)\cr&= {4\over36}+2\left({1\over36}-x^2\right)\cr&= {1\over6}-2x^2. \end{align*}


 Then $2x^2=1/288$, and so $x=1/24$. The probability of obtaining face $F$ is therefore $(1/6)+(1/24)=5/24$, and $m+n=\boxed{29}$.

 




	\item (\textit{AIME II 2006}) Square $ABCD$ has sides of length 1. Points $E$ and $F$ are on $\overline{BC}$ and $\overline{CD},$ respectively, so that $\triangle AEF$ is equilateral. A square with vertex $B$ has sides that are parallel to those of $ABCD$ and a vertex on $\overline{AE}.$ The length of a side of this smaller square is $\frac{a-\sqrt{b}}{c},$ where $a, b,$ and $c$ are positive integers and $b$ is not divisible by the square of any prime. Find $a+b+c.$


 
 \textbf{Answer: }012



\textbf{Solution 1}


                    \begin{figure}[H]
                        \centering
                        \includegraphics[width=0.5\textwidth]{Solution/AIME_II_2006_6_Solution_Figure_1.png}
                    \end{figure}
                    Call the vertices of the new square A', B', C', and D', in relation to the vertices of $ABCD$, and define $s$ to be one of the sides of that square. Since the sides are \href{https://artofproblemsolving.com/wiki/index.php?title=Parallel}{parallel}, by \href{https://artofproblemsolving.com/wiki/index.php?title=Corresponding\_angles&action=edit&redlink=1}{corresponding angles} and AA~ we know that triangles $AA'D'$ and $D'C'E$ are similar. Thus, the sides are proportional: $\frac{AA'}{A'D'} = \frac{D'C'}{C'E} \Longrightarrow \frac{1 - s}{s} = \frac{s}{1 - s - CE}$. Simplifying, we get that $s^2 = (1 - s)(1 - s - CE)$.

 $\angle EAF$ is $60$ degrees, so $\angle BAE = \frac{90 - 60}{2} = 15$. Thus, $\cos 15 = \cos (45 - 30) = \frac{\sqrt{6} + \sqrt{2}}{4} = \frac{1}{AE}$, so $AE = \frac{4}{\sqrt{6} + \sqrt{2}} \cdot \frac{\sqrt{6} - \sqrt{2}}{\sqrt{6} - \sqrt{2}} = \sqrt{6} - \sqrt{2}$. Since $\triangle AEF$ is \href{https://artofproblemsolving.com/wiki/index.php?title=Equilateral}{equilateral}, $EF = AE = \sqrt{6} - \sqrt{2}$. $\triangle CEF$ is a $45-45-90 \triangle$, so $CE = \frac{AE}{\sqrt{2}} = \sqrt{3} - 1$. Substituting back into the equation from the beginning, we get $s^2 = (1 - s)(2 - \sqrt{3} - s)$, so $(3 - \sqrt{3})s = 2 - \sqrt{3}$. Therefore, $s = \frac{2 - \sqrt{3}}{3 - \sqrt{3}} \cdot \frac{3 + \sqrt{3}}{3 + \sqrt{3}} = \frac{3 - \sqrt{3}}{6}$, and $a + b + c = 3 + 3 + 6 = \boxed{12}$.

 Here's an alternative geometric way to calculate $AE$ (as opposed to \href{https://artofproblemsolving.com/wiki/index.php?title=Trigonometry}{trigonometric}): The diagonal $\overline{AC}$ is made of the \href{https://artofproblemsolving.com/wiki/index.php?title=Altitude}{altitude} of the equilateral triangle and the altitude of the $45-45-90 \triangle$. The former is $\frac{AE\sqrt{3}}{2}$, and the latter is $\frac{AE}{2}$; thus $\frac{AE\sqrt{3} + AE}{2} = AC = \sqrt{2} \Longrightarrow AE= \sqrt{6}-\sqrt{2}$. The solution continues as above.

 





\textbf{Solution 2}

Since $\triangle AFE$ is equilateral, $\overline{AE} = \overline{AF}$. It follows that $\overline{FC} = \overline{EC}$. Let $\overline{FC} = x$. Then, $\overline{EF} = x\sqrt{2}$ and $\overline{DF} = 1-x$.

 $\overline{AF} = \sqrt{1+(1-x)^2} = x\sqrt{2}$. 

 Square both sides and combine/move terms to get $x^2+2x-2 = 0$.Therefore $x = -1 + \sqrt{3}$ and $x = -1 - \sqrt{3}$. The second solution is obviously extraneous, so $x = -1 + \sqrt{3}$.

 Now, consider the square ABCD to be on the Cartesian Coordinate Plane with $A = (0,0)$. Then, the line containing $\overline{AF}$ has slope $\frac{1}{2-\sqrt{3}}$ and equation $y = \frac{1}{2-\sqrt{3}}x$.

 The distance from $\overline{DC}$ to $\overline{AF}$ is the distance from $y = 1$ to $y = \frac{1}{2-\sqrt{3}}x$.

 Similarly, the distance from $\overline{AD}$ to $\overline{AF}$ is the distance from $x = 0$ to $y = \frac{1}{2-\sqrt{3}}x$.

 For some value $x = s$, these two distances are equal. 

 $(s-0) = (1 - (\frac{1}{2-\sqrt{3}})s)$

 Solving for s, $s = \frac{3 - \sqrt{3}}{6}$, and $a + b + c = 3 + 3 + 6 = 12$.

 





\textbf{Solution 3}

Suppose $\overline{AB} = \overline{AD} = x.$ Note that $\angle EAF = 60$ since the triangle is equilateral, and by symmetry, $\angle BAE = \angle DAF = 15.$ Note that if $\overline{AD} = x$ and $\angle BAE = 15$, then $\overline{AA'}=\frac{x}{\tan(15)}.$ Also note that 
 \[AB = 1 = \overline{AA'} + \overline{A'B} = \frac{x}{\tan(15)} + x\]Using the fact $\tan(15) = 2-\sqrt{3}$, this yields 
 \[x = \frac{1}{3+\sqrt{3}} = \frac{3-\sqrt{3}}{6} \rightarrow 3 + 3 + 6 = \boxed{12}\]l

 \textbf{Elegant Solution} Why not solve in terms of the side $x$ only (single-variable beauty)? By similar triangles we obtain that $BE=\frac{x}{1-x}$, therefore $CE=\frac{1-2x}{1-x}$. Then $AE=\sqrt{2}*\frac{1-2x}{1-x}$. Using Pythagorean Theorem on $\triangle{ABE}$ yields $\frac{x^2}{(1-x)^2} + 1 = 2 * \frac{(1-2x)^2}{(1-x)^2}$. This means $6x^2-6x+1=0$, and it's clear we take the smaller root: $x=\frac{3-\sqrt{3}}{6}$. Answer: $\boxed{12}$.

 





\textbf{Solution 5 (First part is similar to Solution 2)}

Since $AEF$ is equilateral, $AE=EF$. Let $BE=x$. By the \href{https://artofproblemsolving.com/wiki/index.php?title=Pythagorean\_theorem}{Pythagorean theorem}, $1+x^2=2(1-x)^2$. Simplifying, we get $x^2-4x+1=0$. By the quadratic formula, the roots are $2 \pm \sqrt{3}$. Since $x<1$, we discard the root with the "+", giving $x=2-\sqrt{3}$.
                    \begin{figure}[H]
                        \centering
                        \includegraphics[width=0.5\textwidth]{Solution/AIME_II_2006_6_Solution_Figure_2.png}
                    \end{figure}
                     Let the side length of the square be s. Since $MEK$ is similar to $ABE$, $s=\frac{2-\sqrt{3}-s}{2-\sqrt{3}}$. Solving, we get $s=\frac{3-\sqrt{3}}{6}$ and the final answer is $\boxed{012}$.

 





\textbf{Solution 6 (substitution)}

The equilateral triangle formed is symmetric about the centre of the square. Therefore, $DF = BE = a$ and $FC = CE = b$, so we have $a+b = 1$ and $2b^2 = a^2 +1$. Solving, the only viable solution is $a = 2 -\sqrt{3}$ and $b = \sqrt{3} - 1$.

 Now, let the side length of the second square be $p$ and the remaining side length that composes $a$ be $q$. Then, through similarity, we have $\frac{p}{1-p} = p+q$ and $p+q = 2 - \sqrt{3}$. Substituting, we have $p = (2-\sqrt{3})(1-p)$ and solving gives $p = \frac{3-\sqrt{3}}{6} \implies \boxed{012}$.

 




	\item (\textit{AIME II 2006}) Find the number of ordered pairs of positive integers $(a,b)$ such that $a+b=1000$ and neither $a$ nor $b$ has a zero digit.


 
 \textbf{Answer: }738



\textbf{Solution 1}

There are $\left\lfloor\frac{999}{10}\right\rfloor = 99$ numbers up to 1000 that have 0 as their units digit. All of the other excluded possibilities are when $a$ or $b$ have a 0 in the tens digit, and since the equation is symmetric, we will just count when $a$ has a 0 in the tens digit and multiply by 2 (notice that the only time both $a$ and $b$ can have a 0 in the tens digit is when they are divisible by 100, which falls into the above category, so we do not have to worry about \href{https://artofproblemsolving.com/wiki/index.php?title=Principle\_of\_Inclusion-Exclusion}{overcounting}). 

 Excluding the numbers divisible by 100, which were counted already, there are $9$ numbers in every hundred numbers that have a tens digit of 0 (this is true from 100 to 900), totaling $9 \cdot 9 = 81$ such numbers; considering $b$ also and we have $81 \cdot 2 = 162$. Therefore, there are $999 - (99 + 162) = \boxed{738}$ such ordered pairs.

 





\textbf{Solution 2}

Let $a = \overline{cde}$ and $b = \overline{fgh}$ be 3 digit numbers:

 \begin{verbatim}
cde+fgh----1000
\end{verbatim}$e$ and $h$ must add up to $10$, $d$ and $g$ must add up to $9$, and $c$ and $f$ must add up to $9$. Since none of the digits can be 0, there are $9 \times 8 \times 8=576$ possibilites if both numbers are three digits.

 There are two other scenarios. $a$ and $b$ can be a three digit number and a two digit number, or a three digit number and a one digit number. For the first scenario, there are $9 \times 8 \times 2=144$ possibilities (the two accounting for whether $a$ or $b$ has three digits) and for the second case there are $9 \times 2=18$ possibilities. Thus, thus total possibilities for $(a,b)$ is $576+144+18=738$.

 





\textbf{Solution 3}

We first must notice that we can find all the possible values of $a$ between $1$ and $500$ and then double that result. 

 When $1 < a < 100$ there are $9\times9 = 81$ possible solution for $a$ so that neither $a$ nor $b$ has a zero in it, counting $1$ through $9$, $11$ through $19$, ..., $81$ through $89$.When $100 < a < 200$ there are $9\times8 =72$ possible solution for a so that neither a nor b has a zero in it, counting  $111$ through $119$, $121$ through $129$, ..., $181$ through $189$.This can clearly be extended to $100k < a < 100(k+1)$ where $k$ is an integer and $0 <k < 9$. Thus for $100 < a < 500$ there are $72\times4$ = $288$ possible values of $a$.

 Thus when $1 < a < 500$ there are $288 + 81 =369$ possible values of $a$ and $b$. 

 Doubling this yields $369\times2= 738$.

 





\textbf{Solution 4 (Similar to Solution 2)}

We proceed by casework on the number of digits of $a.$ 

 Case 1: Both $a$ and $b$ have three digits

 We now use constructive counting. For the hundreds digit of $a,$ we see that there are $8$ options - the numbers $1$ through $8.$ (If $a = 9,$ that means that $b$ will be a two digit number, and if $a = 0,$ $a$ will have two digits). Similarly, the tens digit can be $1-8$ as well because a tens digit of $0$ is obviously prohibited and a tens digit of $9$ will lead to a tens digit of $0$ in the other number. The units digit can be anything from $1-9.$ Hence, there are $8 \cdot 8 \cdot 9 = 576$ possible values in this case. 

 Case 2: $a$ (or $b$) has two digits 

 If $a$ has two digits, the only restrictions are that the units digit must not be $0$ and the tens digit must not be $9$ (because then that would lead to $b$ beginning with $90...$). There thus are $8 \cdot 9 = 72$ possibilities for $a,$ and we have to multiply by $2$ because there are the same number of possibilities for $b.$ Thus, there are $72 \cdot 2 = 144$ possible values in this case. 

 Case 3: $a$ (or $b$) has one digit

 This is easy -- $a$ can be anything from $1$ to $9,$ for a total of $9$ possible values. We multiply this by $2$ to account for the single digit $b$ values, so we have $9 \cdot 2 = 18$ possible values for this case. 

 Adding them all up, we get $576 + 144 + 18 = \boxed{738},$ and we're done. 

 Solution by Ilikeapos

 





\textbf{Solution 5}

For every $a \in [1, 999]$, $(a, 1000 - a)$ is a potential candidate for a solution, barring the cases where $a$ or $1000 -a$ has zero digits.

 First, let's consider all viable $a$ that do not have a zero digit. As there are 9 non-zero digits, we have:

 
                \begin{itemize}
                    \item $9^1$ 1-digit numbers without a zero digit\item $9^2$ 2-digit numbers without a zero digit\item $9^3$ 3-digit numbers without a zero digit
                \end{itemize}


                However, we have still overlooked the cases where $1000 - a$ contains zero digits:

 
                \begin{itemize}
                    \item \textbf{Case 1}: If the one's digit of $1000 - a$ is zero, then $a$ will trivially end in a zero, which we've already excluded.\item \textbf{Case 2}: If the ten's digit of $1000 - a$ is zero, with digital representation $\overline{X0Y}$, then $a$ has the digital representation $\overline{[9-X]9[10-Y]}$. $X$ and $Y$ can each take on any value in $[1, 9]$ to produce a value in our set of potential $a$. There are thus $9^2$ cases that we overlooked, where $a$ had no zero digits, but $1000 - a$ did.
                \end{itemize}


                Adding up the cases with $a \in [1, 999]$ with no zero digits and removing the cases with $1000 - a$ with zero digits gives us 
 \[(9^1 + 9^2 + 9^3) - 9^2 = \boxed{738}\]

 





\textbf{Solution 6 (Casework and Constructive Counting, similar to 4)}

Analyze three digit $a$ and $b$ is forced three digits. It follows that the possible units digits would be $(1,9), (2,8), \dots, (9,1)$ which is 9 options total. Then, we carry a 1, so the possible tens digits would be $(1,8), (2,7), \dots, (8,1)$, which is 8 options. We again carry a 1, and the possible hundreds digits is the same as the possible tens digits, so another 8 cases. We have $9 \cdot 8 \cdot 8 = 576$ cases here.

 Now, consider $a$ being a units digit, forcing $b$ as a tens digit, or vice versa. In that case, all integers $1 - 9$ for $a$ work, so we have 9 cases, and for the vice versa case, 18 total. Now, consider $a$ as a two digit forcing $b$ as a three digit, or vice versa. From there, we have that when $a \in {91, 92, \dots, 99}$ then $b$ contains a 0. Otherwise, from $[11, 99]$ inclusive 81 of these values pertain no 0 (through quick constructive counting). From there, that is $81 - 9 = 72$ values, times the vice versa is 144. 

 Thus there are $576+144+18 = \boxed{738}$ ordered pairs.

 




	\item (\textit{AIME II 2006}) There is an unlimited supply of congruent equilateral triangles made of colored paper. Each triangle is a solid color with the same color on both sides of the paper. A large equilateral triangle is constructed from four of these paper triangles. Two large triangles are considered distinguishable if it is not possible to place one on the other, using translations, rotations, and/or reflections, so that their corresponding small triangles are of the same color.


 Given that there are six different colors of triangles from which to choose, how many distinguishable large equilateral triangles may be formed?

                    \begin{figure}[H]
                        \centering
                        \includegraphics[width=0.5\textwidth]{Problem/AIME_II_2006_8_Problem_Figure_1.png}
                    \end{figure}
                    


 
 \textbf{Answer: }336



\textbf{Solution 1}

If two of our big equilateral triangles have the same color for their center triangle and the same \href{https://artofproblemsolving.com/wiki/index.php?title=Multiset}{multiset} of colors for their outer three triangles, we can carry one onto the other by a combination of rotation and reflection.  Thus, to make two triangles distinct, they must differ either in their center triangle or in the collection of colors which make up their outer three triangles.

 There are 6 possible colors for the center triangle.

 
\textbf{Case 1: All three outer triangles have different colors.}

 We can choose the $3$ colors for the outer triangles out of the $6$ total colors in ${6 \choose 3} = 20$ possible choices.

 \textbf{Case 2: Two outer triangles are one color and the third is a different color.}

 We need to choose the two colors for the triangles, then consider a color can be put on two outer triangles or the one outer triangle for $2 {6 \choose 2} = 6 \cdot 5 = 30$ ways.

 \textbf{Case 3: All three outer triangles are the same color.}

 We choose one color for all three outer triangles in ${6\choose 1} = 6$ possible choices.

 
Thus, in total we have $6\cdot(20 + 30 + 6) = \boxed{336}$ total possibilities.

 \textbf{} 




	\item (\textit{AIME II 2006}) Circles $\mathcal{C}_1, \mathcal{C}_2,$ and $\mathcal{C}_3$ have their centers at $(0,0)$, $(12,0)$, and $(24,0)$, and have radii 1, 2, and 4, respectively. Line $t_1$ is a common internal tangent to $\mathcal{C}_1$ and $\mathcal{C}_2$ and has a positive slope, and line $t_2$ is a common internal tangent to $\mathcal{C}_2$ and $\mathcal{C}_3$ and has a negative slope. Given that lines $t_1$ and $t_2$ intersect at $(x,y),$ and that $x=p-q\sqrt{r},$ where $p, q,$ and $r$ are positive integers and $r$ is not divisible by the square of any prime, find $p+q+r.$


 
 \textbf{Answer: }027



\textbf{Solution 1}


                    \begin{figure}[H]
                        \centering
                        \includegraphics[width=0.5\textwidth]{Solution/AIME_II_2006_9_Solution_Figure_1.png}
                    \end{figure}
                    

 Call the centers $O_1, O_2, O_3$, the points of tangency $r_1, r_2, s_1, s_2$ (with $r$ on $t_1$ and $s$ on $t_2$, and $s_2$ on $\mathcal{C}_2$), and the intersection of each common internal tangent to the \href{https://artofproblemsolving.com/wiki/index.php?title=X-axis}{X-axis} $r, s$. $\triangle O_1r_1r \sim \triangle O_2r_2r$ since both triangles have a \href{https://artofproblemsolving.com/wiki/index.php?title=Right\_angle}{right angle} and have vertical angles, and the same goes for $\triangle O_2s_2s \sim \triangle O_3s_1s$. By \href{https://artofproblemsolving.com/wiki/index.php?title=Proportion}{proportionality}, we find that $O_1r = 4$; solving $\triangle O_1r_1r$ by the \href{https://artofproblemsolving.com/wiki/index.php?title=Pythagorean\_theorem}{Pythagorean theorem} yields $r_1r = \sqrt{15}$. On $\mathcal{C}_3$, we can do the same thing to get $O_3s_1 = 4$ and $s_1s = 4\sqrt{3}$.

 The vertical \href{https://artofproblemsolving.com/wiki/index.php?title=Altitude}{altitude} of each of $\triangle O_1r_1r$ and $\triangle O_3s_1s$ can each by found by the formula $c \cdot h = a \cdot b$ (as both products equal twice of the area of the triangle). Thus, the respective heights are $\frac{\sqrt{15}}{4}$ and $2\sqrt{3}$. The horizontal distance from each altitude to the intersection of the tangent with the x-axis can also be determined by the Pythagorean theorem: $\sqrt{15 - \frac{15}{16}} = \frac{15}{4}$, and by 30-60-90: $6$. 

 From this information, the slope of each tangent can be uncovered. The slope of $t_1 = \frac{\Delta y}{\Delta x} = \frac{\frac{\sqrt{15}}{4}}{\frac{15}{4}} = \frac{1}{\sqrt{15}}$. The slope of $t_2 = -\frac{2\sqrt{3}}{6} = -\frac{1}{\sqrt{3}}$.

 The equation of $t_1$ can be found by substituting the point $r (4,0)$ into $y = \frac{1}{\sqrt{15}}x + b$, so $y = \frac{1}{\sqrt{15}}x - \frac{4}{\sqrt{15}}$. The equation of $t_2$, found by substituting point $s (16,0)$, is $y = \frac{-1}{\sqrt{3}}x + \frac{16}{\sqrt{3}}$. Putting these two equations together results in the desired $\frac{1}{\sqrt{15}}x - \frac{4}{\sqrt{15}} = -\frac{1}{\sqrt{3}}x + \frac{16}{\sqrt{3}}$ $\Longrightarrow x = \frac{16\sqrt{5} + 4}{\sqrt{5} + 1} \cdot \frac{\sqrt{5}-1}{\sqrt{5}-1}$ $= \frac{76 - 12\sqrt{5}}{4}$ $= 19 - 3\sqrt{5}$. Thus, $p + q + r = 19 + 3 + 5 \Longrightarrow \boxed{027}$.

 





\textbf{Solution 2}


                    \begin{figure}[H]
                        \centering
                        \includegraphics[width=0.5\textwidth]{Solution/AIME_II_2006_9_Solution_Figure_1.png}
                    \end{figure}
                    

 This solution uses the diagram of solution 1

 $t_1$ and $t_2$ must traverse circle $\mathcal{C}_1$ and $\mathcal{C}_2$, and $\mathcal{C}_2$ and $\mathcal{C}_3$ internal points of similitude, say $R$ and $S$. As the circles lie with centre's on the $x$-axis, it is obvious that $R$ must have a $y$ coordinate 0 and so must $S$. Also, the distance from the centre of $\mathcal{C}_2$ to that of $R$ and the distance from $S$ to $\mathcal{C}_1$ must be the ratio of their radii due to definition of points of similitude and the same applies with point $S$. Then, one can find that $\frac{R_x}{12-R_x} = 0.5$, so $R_x = 4$, and applying the same for $S_x = 16$. Therefore, the internal points of similitude for the pairwise circles $\mathcal{C}_1, \mathcal{C}_2$, and $\mathcal{C}_2, \mathcal{C}_3$ is $R(4,0), S(16,0)$.

 We then attempt to find the slope of $t_1$ and $t_2$, knowing they go through $R$ and $S$. We can quickly see that the right triangle formed by the tangential point of the radius of $\mathcal{C}_3$ and $t_2$, the centre of $\mathcal{C}_3$, and $S$ is a 30-60-90 degree triangle. The slope of this line is $\tan(60^\circ)$ (positive in the third quadrant), which is $\sqrt{3}$, and so the slope of $t_2$ is $-\frac{1}{\sqrt{3}}$. We know it passes through $S(16,0)$, so the equation of $t_2$ is $y = -\frac{1}{\sqrt{3}}x + \frac{16}{\sqrt{3}}$. Now, $t_1$ is where things get slightly trickier. Luckily, all we have to do is find $-\tan(RO_1R_1)$, which is just $-\sqrt{15}$. It follows that the slope of $t_1$ is just $\frac{1}{\sqrt{15}}$, and it passes through $R(4,0)$, so the equation of line $t_1$ is $y = \frac{1}{\sqrt{15}}x - \frac{4}{\sqrt{15}}$. Solving for the intersection $(x,y)$, the solution is $(x,y) = (19-3\sqrt{5}, \sqrt{15}- \sqrt{3})$, so $x = 19 - 3\sqrt{5}$. The requested value is $19+3+5 = \boxed{027}$.

 




	\item (\textit{AIME II 2006}) Seven teams play a soccer tournament in which each team plays every other team exactly once. No ties occur, each team has a $50\%$ chance of winning each game it plays, and the outcomes of the games are independent. In each game, the winner is awarded a point and the loser gets 0 points. The total points are accumulated to decide the ranks of the teams. In the first game of the tournament, team $A$ beats team $B.$ The probability that team $A$ finishes with more points than team $B$ is $m/n,$ where $m$ and $n$ are relatively prime positive integers. Find $m+n.$


 
 \textbf{Answer: }831



\textbf{Solution 1}

The results of the five remaining games are independent of the first game, so by symmetry, the probability that $A$ scores higher than $B$ in these five games is equal to the probability that $B$ scores higher than $A$. We let this probability be $p$; then the probability that $A$ and $B$ end with the same score in these five games is $1-2p$.

 Of these three cases ($|A| > |B|, |A| < |B|, |A|=|B|$), the last is the easiest to calculate (see solution 2 for a way to directly calculate the other cases). 

 There are ${5\choose k}$ ways to $A$ to have $k$ victories, and ${5\choose k}$ ways for $B$ to have $k$ victories. Summing for all values of $k$,

 $1-2p = \frac{1}{2^{5} \times 2^{5}}\left(\sum_{k=0}^{5} {5\choose k}^2\right) = \frac{1^2+5^2+10^2+10^2+5^2+1^2}{1024} = \frac{126}{512}.$Thus $p = \frac 12 \left(1-\frac{126}{512}\right) = \frac{193}{512}$. The desired probability is the sum of the cases when $|A| \ge |B|$, so the answer is $\frac{126}{512} + \frac{193}{512} = \frac{319}{512}$, and $m+n = \boxed{831}$.

 





\textbf{Solution 2}

You can break this into cases based on how many rounds $A$ wins out of the remaining $5$ games.

 
                \begin{itemize}
                    \item If $A$ wins 0 games, then $B$ must win 0 games and the probability of this is $\frac{{5 \choose 0}}{2^5} \frac{{5 \choose 0}}{2^5} = \frac{1}{1024}$.
                \end{itemize}


                
                \begin{itemize}
                    \item If $A$ wins 1 games, then $B$ must win 1 or less games and the probability of this is $\frac{{5 \choose 1}}{2^5} \frac{{5 \choose 0}+{5 \choose 1}}{2^5} = \frac{30}{1024}$.
                \end{itemize}


                
                \begin{itemize}
                    \item If $A$ wins 2 games, then $B$ must win 2 or less games and the probability of this is $\frac{{5 \choose 2}}{2^5} \frac{{5 \choose 0}+{5 \choose 1}+{5 \choose 2}}{2^5} = \frac{160}{1024}$.
                \end{itemize}


                
                \begin{itemize}
                    \item If $A$ wins 3 games, then $B$ must win 3 or less games and the probability of this is $\frac{{5 \choose 3}}{2^5} \frac{{5 \choose 0}+{5 \choose 1}+{5 \choose 2}+{5 \choose 3}}{2^5} = \frac{260}{1024}$.
                \end{itemize}


                
                \begin{itemize}
                    \item If $A$ wins 4 games, then $B$ must win 4 or less games and the probability of this is $\frac{{5 \choose 4}}{2^5} \frac{{5 \choose 0}+{5 \choose 1}+{5 \choose 2}+{5 \choose 3}+{5 \choose 4}}{2^5} = \frac{155}{1024}$.
                \end{itemize}


                
                \begin{itemize}
                    \item If $A$ wins 5 games, then $B$ must win 5 or less games and the probability of this is $\frac{{5 \choose 5}}{2^5} \frac{{5 \choose 0}+{5 \choose 1}+{5 \choose 2}+{5 \choose 3}+{5 \choose 4}+{5 \choose 5}}{2^5} = \frac{32}{1024}$.
                \end{itemize}


                Summing these 6 cases, we get $\frac{638}{1024}$, which simplifies to $\frac{319}{512}$, so our answer is $319 + 512 = \boxed{831}$.

 





\textbf{Solution 3}

We can apply the concept of generating functions here. 

 The generating function for $B$ is $(1 + 0x^{1})$ for the first game where $x^{n}$ is winning n games. Since $B$ lost the first game, the coefficient for $x^{1}$ is 0. The generating function for the next 5 games is $(1 + x)^{5}$. Thus, the total generating function for number of games he wins is

 \begin{verbatim}
${(1 + 0x)(1 + x)^{5}} = (1 + 5x^{1} + 10x^{2} + 10x^{3} + 5x^{4} + x^{5})$.
\end{verbatim}The generating function for $A$ is the same except that it is multiplied by $x$ instead of $(1+0x)$. Thus, the generating function for $A$ is 

 $1x + 5x^{2} + 10x^{3} + 10x^{4} + 5x^{5} + x^{6}$. 

 The probability that $B$ wins 0 games is $\frac{1}{32}$. Since the coefficients for all $x^{n}$ where 

 $n \geq 1$ sums to 32, the probability that $A$ wins more games is $\frac{32}{32}$. 

 Thus, the probability that $A$ has more wins than $B$ is $\frac{1}{32} \times \frac{32}{32} + \frac{5}{32} \times \frac{31}{32} + \frac{10}{32} \times \frac{26}{32} + \frac{10}{32} \times \frac{16}{32} + \frac{5}{32} \times \frac{6}{32} +\frac{1}{32} \times \frac{1}{32} = \frac{638}{1024} = \frac{319}{512}$.

 Thus, $319 + 512 = \boxed{831}$.

 





\textbf{Solution 4}

After the first game, there are $10$ games we care about-- those involving $A$ or $B$. There are $3$ cases of these $10$ games: $A$ wins more than $B$, $B$ wins more than $A$, or $A$ and $B$ win the same number of games. Also, there are $2^{10} = 1024$ total outcomes. By symmetry, the first and second cases are equally likely, and the third case occurs $\binom{5}{0}^2+\binom{5}{1}^2+\binom{5}{2}^2+\binom{5}{3}^2+\binom{5}{4}^2+\binom{5}{5}^2 = \binom{10}{5} = 252$ times, by \href{https://artofproblemsolving.com/wiki/index.php?title=Combinatorial\_identity\#Another\_Identity}{a special case of Vandermonde's Identity}. There are therefore $\frac{1024-252}{2} = 386$ possibilities for each of the other two cases.

 If $B$ has more wins than $A$ in its $5$ remaining games, then $A$ cannot beat $B$ overall. However, if $A$ has more wins or if $A$ and $B$ are tied, $A$ will beat $B$ overall. Therefore, out of the $1024$ possibilites, $386+252 = 638$ ways where $A$ wins, so the desired probability is $\frac{638}{1024} = \frac{319}{512}$, and $m+n=\boxed{831}$.

 




	\item (\textit{AIME II 2006}) A sequence is defined as follows $a_1=a_2=a_3=1,$ and, for all positive integers $n, a_{n+3}=a_{n+2}+a_{n+1}+a_n.$ Given that $a_{28}=6090307, a_{29}=11201821,$ and $a_{30}=20603361,$ find the remainder when $\sum^{28}_{k=1} a_k$ is divided by 1000.


 
 \textbf{Answer: }834



\textbf{Solution 1}

Define the sum as $s$. Since $a_n\ = a_{n + 3} - a_{n + 2} - a_{n + 1}$, the sum will be:

 $s = a_{28} + \sum^{27}_{k=1} (a_{k+3}-a_{k+2}-a_{k+1}) \\ s = a_{28} + \left(\sum^{30}_{k=4} a_{k} - \sum^{29}_{k=3} a_{k}\right) - \left(\sum^{28}_{k=2} a_{k}\right)\\ s = a_{28} + (a_{30} - a_{3}) - \left(\sum^{28}_{k=2} a_{k}\right) = a_{28} + a_{30} - a_{3} - (s - a_{1})\\ s = -s + a_{28} + a_{30}$Thus $s = \frac{a_{28} + a_{30}}{2}$, and $a_{28},\,a_{30}$ are both given; the last four digits of their sum is $3668$, and half of that is $1834$. Therefore, the answer is $\boxed{834}$.−	

 





\textbf{Solution 2 (bash)}

Since the problem only asks for the first 28 terms and we only need to calculate mod 1000, we simply bash the first 28 terms:

 


 $a_{1}\equiv 1 \pmod {1000} \\ a_{2}\equiv 1 \pmod {1000} \\ a_{3}\equiv 1 \pmod {1000} \\ a_{4}\equiv 3 \pmod {1000} \\ a_{5}\equiv 5 \pmod {1000} \\ a_{6}\equiv 9 \pmod {1000} \\ a_{7}\equiv 17 \pmod {1000} \\ a_{8}\equiv 31 \pmod {1000} \\ a_{9}\equiv 57 \pmod {1000} \\ a_{10}\equiv 105 \pmod {1000} \\ a_{11}\equiv 193 \pmod {1000} \\ a_{12}\equiv 355 \pmod {1000} \\ a_{13}\equiv 653 \pmod {1000} \\ a_{14}\equiv 201 \pmod {1000} \\ a_{15}\equiv 209 \pmod {1000} \\ a_{16}\equiv 63 \pmod {1000} \\ a_{17}\equiv 473 \pmod {1000} \\ a_{18}\equiv 745 \pmod {1000} \\ a_{19}\equiv 281 \pmod {1000} \\ a_{20}\equiv 499 \pmod {1000} \\ a_{21}\equiv 525 \pmod {1000} \\ a_{22}\equiv 305 \pmod {1000} \\ a_{23}\equiv 329 \pmod {1000} \\ a_{24}\equiv 159 \pmod {1000} \\ a_{25} \equiv 793 \pmod {1000} \\ a_{26} \equiv 281 \pmod {1000} \\ a_{27} \equiv 233 \pmod {1000} \\ a_{28} \equiv 307 \pmod {1000}$

 Adding all the residues shows the sum is congruent to $\boxed{834}$ mod $1000$.

 





\textbf{Solution 3 (some guessing involved)/"Engineer's Induction"}

All terms in the sequence are sums of previous terms, so the sum of all terms up to a certain point must be some linear combination of the first three terms. Also, we are given $a_{28}, a_{29},$ and $a_{30}$, so we can guess that there is some way to use them in a formula. Namely, we guess that there exists some $p, q, r$ such that $\sum_{k=1}^{n}{a_k} = pa_n+qa_{n+1}+ra_{n+2}$. From here, we list out the first few terms of the sequence and the cumulative sums, and with a little bit of substitution and algebra we see that $(p, q, r) = (\frac{1}{2}, 0, \frac{1}{2})$, at least for the first few terms. From this, we have that $\sum_{k=1}^{28}{a_k} = \frac{a_{28}+a_{30}}{2} \equiv{\boxed{834}}\pmod {1000}$.

 Solution by zeroman; clarified by srisainandan6

 


 





\textbf{Solution 4 (if you did not know how to use the numbers given in the problem)}

By the Chinese remainder theorem, each number under 1000 is uniquely determined by its mod 8 and mod 125.

 We list a few terms of the sequence mod 8:

 
 \[1,1,1,3,5,1,1,7,1,1,1,...\]

 Therefore, the cycle repeats every 8 numbers, and each cycle has a sum of 4 mod 8. Therefore, the sum mod 8 is 
 \[3 \cdot 4 + 1 + 1 + 1 + 3 = 2 \mod 8.\]

 Denote 
 \[s_n = \sum _{i=1} ^{n} a_i.\] It is easy to prove that $s_{n+3} = s_{n+2} + s_{n+1} + s_{n}.$

 We write the sum ($s_n$) of the first terms mod 125:

 
 \[1,2,3,6,11,20,37,68,0,-20,48,28,56,7,91,29,2,-3,28,27,52,-18,61,-30,13,44,27,84.\]

 Therefore the desired number is $84 \mod 125.$

 From here, we can determine the number we are looking for is $750 + 84 = \boxed{834}.$-sd8

 




	\item (\textit{AIME II 2006}) Equilateral $\triangle ABC$ is inscribed in a circle of radius 2. Extend $\overline{AB}$ through $B$ to point $D$ so that $AD=13,$ and extend $\overline{AC}$ through $C$ to point $E$ so that $AE = 11.$ Through $D,$ draw a line $l_1$ parallel to $\overline{AE},$ and through $E,$ draw a line $l_2$ parallel to $\overline{AD}.$ Let $F$ be the intersection of $l_1$ and $l_2.$ Let $G$ be the point on the circle that is collinear with $A$ and $F$ and distinct from $A.$ Given that the area of $\triangle CBG$ can be expressed in the form $\frac{p\sqrt{q}}{r},$  where $p, q,$ and $r$ are positive integers, $p$ and $r$  are relatively prime, and $q$ is not divisible by the square of any prime, find $p+q+r.$


 
 \textbf{Answer: }865



\textbf{Solution 1}


                    \begin{figure}[H]
                        \centering
                        \includegraphics[width=0.5\textwidth]{Solution/AIME_II_2006_12_Solution_Figure_1.png}
                    \end{figure}
                    Notice that $\angle{E} = \angle{BGC} = 120^\circ$ because $\angle{A} = 60^\circ$. Also, $\angle{GBC} = \angle{GAC} = \angle{FAE}$ because they both correspond to arc ${GC}$. So $\Delta{GBC} \sim \Delta{EAF}$.

 
 \[[EAF] = \frac12 (AE)(EF)\sin \angle AEF  = \frac12\cdot11\cdot13\cdot\sin{120^\circ} = \frac {143\sqrt3}4.\]

 Because the ratio of the area of two similar figures is the square of the ratio of the corresponding sides, $[GBC] = \frac {BC^2}{AF^2}\cdot[EAF] = \frac {12}{11^2 + 13^2 - 2\cdot11\cdot13\cdot\cos120^\circ}\cdot\frac {143\sqrt3}4 = \frac {429\sqrt3}{433}$. Therefore, the answer is $429+433+3=\boxed{865}$.

 





\textbf{Solution 2: Analytic Geometry/Coord Bash}

Solution by e\_power\_pi\_times\_i/edited by srisainandan6

 Let the center of the circle be $O$ and the origin. Then, $A (0,2)$, $B (-\sqrt{3}, -1)$, $C (\sqrt{3}, -1)$. $D$ and $E$ can be calculated easily knowing $AD$ and $AE$, $D (-\dfrac{13}{2}, \dfrac{-13\sqrt{3}+4}{2})$, $E (\dfrac{11}{2}, \dfrac{-11\sqrt{3}+4}{2})$. As $DF$ and $EF$ are parallel to $AE$ and $AD$, $F (-1, -12\sqrt{3}+2)$. $G$ and $A$ is the intersection between $AF$ and circle $O$. Therefore $G (-\dfrac{48\sqrt{3}}{433}, -\dfrac{862}{433})$. Using the Shoelace Theorem, $[CBG] = \dfrac{429\sqrt{3}}{433}$, so the answer is $\boxed{865}$. Note that although the solution may appear short, actually getting all the coordinates take a while as there is plenty of computation.

 Note by chrisdiamond10: We can save time calculating the area of the triangle once we have the coordinates of $B,C,G$ by using $\frac{b\cdot h}{2}$. Use $BC$ as the base, then the base is $2\sqrt{3}$. The height is easily calculated as $-1-\left(-\frac{862}{433}\right)=-1+\frac{862}{433}=\frac{429}{433}$, so multiplying base by height and dividing by two we find that the total area is $\frac{429\sqrt{3}}{433}$, and our answer is $\boxed{865}$.

 





\textbf{Solution 3: Trig}

Lines $l_1$ and $l_2$ are constructed such that $AEFD$ is a parallelogram, hence $DF = 13$. Since $BAC$ is equilateral with angle of $60^{\circ}$, angle $D$ is $120^{\circ}$. Use law of cosines to find $AF = \sqrt{433}$. Then use law of sines to find angle $BAG$ and $GAC$. Next we use Ptolemy's Theorem on $ABGC$ to find that $CG + BG = AG$. Next we use law of cosine on triangles $BAG$ and $GAC$, solving for BG and CG respectively. Subtract the two equations and divide out a $BG + CG$ to find the value of $CG - BG$. Next, $AG = 2\cdot R \cos{\theta}$, where R is radius of circle $= 2$ and $\theta =$ angle $BAG$. We already know sine of the angle so find cosine, hence we have found $AG$. At this point it is system of equation yielding $CG = \frac{26\sqrt{3}}{\sqrt{433}}$ and $BG = \frac{22\sqrt{3}}{\sqrt{433}}$. Given $[CBG] = \frac{BC \cdot CG \cdot BG}{4R}$, and $BC = 2\sqrt{3}$ by $30-60-90$ triangle, we can evaluate to find $[CBG] = \frac{429\sqrt{3}}{433}$, to give answer = $\boxed{865}$.

 





\textbf{Solution 4}

Note that $AB=2\sqrt3$, $DF=11$, and $EF=13$. If we take a homothety of the parallelogram with respect to $A$, such that $F$ maps to $G$, we see that $\frac{[ABG]}{[ACG]}=\frac{11}{13}$. Since $\angle AGB=\angle AGC=60^{\circ}$, from the sine area formula we have $\frac{BG}{CG}=\frac{11}{13}$. Let $BG=11k$ and $CG=13k$.

 By Law of Cosines on $\triangle BGC$, we have
 \[12=k^2(11^2+11\cdot13+13^2)=433k^2\implies k^2=\frac{12}{433}\]Thus, $[CBG]=\frac12 (11k)(13k)\sin 120^{\circ} = \frac{\sqrt3}{4}\cdot 143\cdot \frac{12}{433}=\frac{429\sqrt3}{433}\implies\boxed{865}$.

 





\textbf{Solution 5: Easy Similar triangles}

First, by properties of cyclic quadrilterals, $<BGC = 180-<BAC=120$. Next, we notice quadrilateral AEFD is a parallelogram, and opposites angles are supplementary, and therefore $<ADF = 180-<BAC = 120$. Finally, by inscribed angles, $<BAG = <BCG$. Thus we can conclude $\triangle CBG \sim \triangle AFD$ by AA similarity. 

 Because $\triangle FAE \cong \triangle AFD$, $\triangle FAE \sim \triangle CBG$ aswell. So, we can write $\frac{BC}{AF} = \frac{BG}{DF} \to \frac{2\sqrt{3}}{\sqrt{433}} = \frac{BG}{11}$ since $AF = \sqrt{433}$ by the Law of Cosines. This gives $BG = \frac{22\sqrt{3}}{\sqrt{433}}$ and doing an identical process yields $CG = \frac{26\sqrt{3}}{\sqrt{433}}$. Thus, $[CBG] = \frac{26\sqrt{3}}{\sqrt{433}} \cdot \frac{22\sqrt{3}}{\sqrt{433}} \cdot \sin 120 = \frac{429\sqrt{3}}{433} \to p+q+r = \boxed{865}$.

 




	\item (\textit{AIME II 2006}) How many integers $N$ less than 1000 can be written as the sum of $j$ consecutive positive odd integers from exactly 5 values of $j\ge 1$?


 
 \textbf{Answer: }015



\textbf{Solution 1}

Let the first odd integer be $2n+1$, $n\geq 0$. Then the final odd integer is $2n+1 + 2(j-1) = 2(n+j) - 1$. The odd integers form an \href{https://artofproblemsolving.com/wiki/index.php?title=Arithmetic\_sequence}{arithmetic sequence} with sum $N = j\left(\frac{(2n+1) + (2(n+j)-1)}{2}\right) = j(2n+j)$. Thus, $j$ is a factor of $N$.

 Since $n\geq 0$, it follows that $2n+j \geq j$ and $j\leq \sqrt{N}$.

 Since there are exactly $5$ values of $j$ that satisfy the equation, there must be either $9$ or $10$ factors of $N$. This means $N=p_1^2p_2^2$ or $N=p_1p_2^4$. Unfortunately, we cannot simply observe prime factorizations of $N$ because the factor $(2n+j)$ does not cover all integers for any given value of $j$.

 Instead we do some casework:

 
                \begin{itemize}
                    \item If $N$ is odd, then $j$ must also be odd. For every odd value of $j$, $2n+j$ is also odd, making this case valid for all odd $j$. Looking at the forms above and the bound of $1000$, $N$ must be
                \end{itemize}


                
 \[(3^2\cdot5^2),\ (3^2\cdot7^2),\ (3^4\cdot5),\ (3^4\cdot7),\ (3^4\cdot 11)\]

 
            \begin{enumerate}[label=(\roman*)]
                \item Those give $5$ possibilities for odd $N$.
            \end{enumerate}
            
                \begin{itemize}
                    \item If $N$ is even, then $j$ must also be even. Substituting $j=2k$, we get
                \end{itemize}


                
 \[N = 4k(n+k) \Longrightarrow \frac{N}{4} = k(n+k)\]

 
            \begin{enumerate}[label=(\roman*)]
                \item Now we can just look at all the \href{https://artofproblemsolving.com/wiki/index.php?title=Prime\_factorization}{prime factorizations} since $(n+k)$ cover the integers for any $k$. Note that our upper bound is now $250$:
            \end{enumerate}
            
 \[\frac{N}{4} = (2^2\cdot3^2),(2^2\cdot5^2),(2^2\cdot7^2), (3^2\cdot5^2), (2^4\cdot3), (2^4\cdot5), (2^4\cdot7), (2^4\cdot11), (2^4\cdot13), (3^4\cdot2)\]

 
            \begin{enumerate}[label=(\roman*)]
                \item Those give $10$ possibilities for even $N$.
            \end{enumerate}
            The total number of integers $N$ is $5 + 10 = \boxed{15}$.

 





\textbf{Solution 2}

Let the largest odd number below the sequence be the $q$th positive odd number, and the largest odd number in the sequence be the $p$th positive odd number. Therefore, the sum is $p^2-q^2=(p+q)(p-q)$ by sum of consecutive odd numbers. Note that $p+q$ and $p-q$ have the same parity, and $q$ can equal $0$. We then perform casework based on the parity of $p-q$. 

 If $p-q$ is odd, then $p^2-q^2$ must always be odd.  Therefore, to have 5 pairs of odd factors, we must have either $9$ (in which case the number is a perfect square) or $10$ factors. Considering the upper bound, the only way this can happen is $p_1^4\cdot{p_2}$ or $p_1^2\cdot{p_2^2}$. N must then be one of
 \[(3^2\cdot5^2),\ (3^2\cdot7^2),\ (3^4\cdot5),\ (3^4\cdot7),\ (3^4\cdot 11)\]So, there are $5$ solutions when $(p+q)(p-q)$ is odd. 

 If $p-q$ is even, then $(p+q)(p-q)$ must have at least two factors of $2$, so we can rewrite the expression as $4(k)(k-q)$ where $k=\frac{p+q}{2}$. We can disregard the $4$ by dividing by $4$ and restricting our upper bound to $250$. Since $k$ and $q$ don't have to be the same parity, we can include all cases less than 250 that have 9 or 10 factors. We then have 
 \[(2^2\cdot3^2),(2^2\cdot5^2),(2^2\cdot7^2), (3^2\cdot5^2), (2^4\cdot3), (2^4\cdot5), (2^4\cdot7), (2^4\cdot11), (2^4\cdot13), (3^4\cdot2)\]as the possibilities. 

 Therefore, there are $10+5=\boxed{015}$ possibilities for $p^2-q^2=N$ ~sigma

 




	\item (\textit{AIME II 2006}) Let $S_n$ be the sum of the reciprocals of the non-zero digits of the integers from $1$ to $10^n$ inclusive. Find the smallest positive integer $n$ for which $S_n$ is an integer.


 
 \textbf{Answer: }063



\textbf{Solution}

Let $K = \sum_{i=1}^{9}{\frac{1}{i}}$.  Examining the terms in $S_1$, we see that $S_1 = K + 1$ since each digit $n$ appears once and 1 appears an extra time.  Now consider writing out $S_2$.  Each term of $K$ will appear 10 times in the units place and 10 times in the tens place (plus one extra 1 will appear), so $S_2 = 20K + 1$.

 In general, we will have that

 $S_n = (n10^{n-1})K + 1$because each digit will appear $10^{n - 1}$ times in each place in the numbers $1, 2, \ldots, 10^{n} - 1$, and there are $n$ total places.  

 The denominator of $K$ is $D = 2^3\cdot 3^2\cdot 5\cdot 7$.  For $S_n$ to be an integer, $n10^{n-1}$ must be divisible by $D$.  Since $10^{n-1}$ only contains the factors $2$ and $5$ (but will contain enough of them when $n \geq 3$), we must choose $n$ to be \href{https://artofproblemsolving.com/wiki/index.php?title=Divisible}{divisible} by $3^2\cdot 7$.  Since we're looking for the smallest such $n$, the answer is $\boxed{063}$.

 





\textbf{Solution 2 (easier approach but nearly same logic)}

Consider \( n \ge 2 \), since clearly \( n = 1 \) doesn't work.

 For a certain \( 1/i \) to be summed such that it is an integer, we must have \( i \) to be summed continuously \( mi \) times for some integer \( m \). For example, if we have \( 1/5 \), we sum \( 1/5 \) either \( 5, 10, 15, \dots \) times to obtain an integer.

 Thus, every digit \( i \) must be \( 1, 2, 3, \dots, 9 \). It is known that every digit between every period \( 10^k \) appears the same number of times. For example, for the first \( 100 \) numbers (excluding \( 0 \) and \( 100 \)), every digit from \( 1 \text{--} 9 \) appears the same number of times.

 If we can find an algorithm that calculates the number of times a digit appears between \( 1 \) and \( 10^n \), and extend that for all values of \( n \), we could look for a specific value of \( n \) such that every digit \( 1, 2, 3, \dots, 9 \) appears a number of times such that all the reciprocals, when summed, are integers.

 To find this algorithm, we construct a pattern. Consider \( 1 \text{--} 100 \). Let's try to find the number of \( 9 \)'s between this interval. There are \( 9 \) numbers (excluding from \( 90 \text{--} 99 \)), whose units digits is \( 9 \). From \( 90 \text{--} 98 \), we have nine \( 9 \)'s. \( 99 \) gives us \( 2 \) more \( 9 \)'s, and thus we have a total of \( 11 + 9 = 20 \) nines. We then consider \( 1 \text{--} 1000 \). Between \( 1 \) and \( 100 \), we have \( 21 \) nines. The same applies from \( 101 \text{--} 200, 201 \text{--} 300, \dots, 801 \text{--} 899 \). We then have a total of \( 21 \times 9 = 189 \) nines. Then, we have \( 1 \) nine from \( 900 \), so \( 190 \) nines. Then, from \( 901 \text{--} 908, 911 \text{--} 918, 921 \text{--} 928, \dots, 981 \text{--} 988 \), we have one \( 9 \), contributing \( 8 \times 9 = 72 \) nines. For \( 10 \) additional numbers (\( 909, 919, 929, \dots, 989, 990 \)), we have \( 2 \) nines, contributing \( 20 \) nines. Then, for the final set of \( 9 \), we also have \( 2 \) nines each, giving \( 18 \) more nines, and the last \( 999 \) has \( 3 \) nines, giving a total of \( 190 + 72 + 20 + 18 = 300 \) nines.

 Now we must find a pattern. For every number \( 1 \text{--} 10^n \), we must find the number of \( 9 \)'s. Note that for all numbers less than \( 10^n \), we are given that \( 9 \) can appear in any one of the positions between the units digit, and the \( 10^n \) decimal digit. That is \( n \) places. For each \( 10^n \) places, the numbers appear equally often, and thus there are \( 10^n / 10 = 10^{n-1} \) slots. Our algorithm is then \( n \cdot 10^{n-1} \). 

 Note that this formula also works for \( 2, 3, 4, \dots, 8 \). For \( 1 \), note that we have to add \( 1 \) that appears from \( 10^n \), so the formula for \( 1 \) is slightly different, being \( n \cdot 10^{n-1} + 1 \).

 Thus, we now are at the second part. We require \( n \) such that 

 
 \[1 \mid (n \cdot 10^{n-1} + 1) \quad \text{and} \quad 2, 3, 4, 5, \dots, 9 \mid n \cdot 10^{n-1}\]

 Obviously, the former is always true, as \( 1 \) divides everything. Thus, we neglect the former and pay attention to the latter.

 The \( \text{lcm}(2,3,4,5,\dots,9) \), after some work, is 2520. Thus, \( 2520 \mid n \cdot 10^{n-1} \). The prime factorization of 2520 is \( 2^3 \cdot 3^2 \cdot 5 \cdot 7 \). The prime factorization of \( 10^{n-1} \) is \( 2^{n-1} \cdot 5^{n-1} \). To check extraneous, note that for \( n \ge 4 \), we have that \( 2^{n-1} \ge 2^3 \), and the same for \( 5^{n-1} \) with \( n \ge 2 \), hence this is solvable. Because 9 and 7 are not factors of 10, we have that 9 and 7 must divide into \( n \), meaning that \( 63 \mid n \). The smallest \( n \) is then right in front of us, or \( \boxed{063} \).

 Awesome solution but just a clarification - your manual counting for the number of nines from 1-1000 is wrong as you contradicted yourself in the go: here's the correct counting:For 1-899 as you said 190, it is not 190 it is 181 as each interval 1-100 has 20 nines as you said not 21. Hence 20×9=180 and adding a 9 from 900-> 181 9's in total. Now for 901-999:If we analyse be decimal places, each hundred's digit has a 9 so, 99 9s from there and for ten's digit-> obviously the numbers from 990 to 999 only satisfy this criteria so 10 9's from hereAnd for unit's digit, the numbers follow the pattern 909,919,929 and so on till 999 and we get 10 9s from here as well, now if we add the total nines from 901-999 it is 99+10+10=119 And in total, the number of nines from 1-1000 is 181+119=300

 \begin{verbatim}
																							~Randommathdoerfromindia
\end{verbatim}




	\item (\textit{AIME II 2006}) Given that $x, y,$ and $z$ are real numbers that satisfy:

 \begin{align*} x &= \sqrt{y^2-\frac{1}{16}}+\sqrt{z^2-\frac{1}{16}}, \\ y &= \sqrt{z^2-\frac{1}{25}}+\sqrt{x^2-\frac{1}{25}}, \\ z &= \sqrt{x^2 - \frac 1{36}}+\sqrt{y^2-\frac 1{36}}, \end{align*}

and that $x+y+z = \frac{m}{\sqrt{n}},$ where $m$ and $n$ are positive integers and $n$ is not divisible by the square of any prime, find $m+n.$


 
 \textbf{Answer: }009



\textbf{Solution 1 (Geometric Interpretation)}

Let $\triangle XYZ$ be a triangle with sides of length $x, y$ and $z$, and suppose this triangle is acute (so all \href{https://artofproblemsolving.com/wiki/index.php?title=Altitude}{altitudes} are in the interior of the triangle).

 
Let the altitude to the side of length $x$ be of length $h_x$, and similarly for $y$ and $z$. Then we have by two applications of the \href{https://artofproblemsolving.com/wiki/index.php?title=Pythagorean\_Theorem}{Pythagorean Theorem} we that 
 \[x = \sqrt{y^2 - h_x^2} + \sqrt{z^2 - h_x^2}\] As a \href{https://artofproblemsolving.com/wiki/index.php?title=Function}{function} of $h_x$, the \href{https://artofproblemsolving.com/wiki/index.php?title=RHS}{RHS} of this equation is strictly decreasing, so it takes each value in its \href{https://artofproblemsolving.com/wiki/index.php?title=Range}{range} exactly once.  Thus we must have that $h_x^2 = \frac1{16}$ and so $h_x = \frac{1}4$ and similarly $h_y = \frac15$ and $h_z = \frac16$.

 
The area of the triangle must be the same no matter how we measure it; therefore $x\cdot h_x = y\cdot h_y = z \cdot h_z$ gives us $\frac x4 = \frac y5 = \frac z6 = 2A$ and $x = 8A, y = 10A$ and $z = 12A$.  

 
The \href{https://artofproblemsolving.com/wiki/index.php?title=Semiperimeter}{semiperimeter} of the triangle is $s = \frac{8A + 10A + 12A}{2} = 15A$ so by \href{https://artofproblemsolving.com/wiki/index.php?title=Heron\%27s\_formula}{Heron's formula} we have 
 \[A = \sqrt{15A \cdot 7A \cdot 5A \cdot 3A} = 15A^2\sqrt{7}\] 

 Thus, $A = \frac{1}{15\sqrt{7}}$ and $x + y + z = 30A = \frac2{\sqrt{7}}$ and the answer is $2 + 7 = \boxed{009}$.

 The justification that there is an acute triangle with sides of length $x, y$ and $z$:

 Note that $x, y$ and $z$ are each the sum of two positive \href{https://artofproblemsolving.com/wiki/index.php?title=Square\_root}{square roots} of real numbers, so $x, y, z \geq 0$.  (Recall that, by \href{https://artofproblemsolving.com/wiki/index.php?title=AIME}{AIME} \href{https://artofproblemsolving.com/wiki/index.php?title=Mathematical\_convention}{ convention}, all numbers (including square roots) are taken to be real unless otherwise indicated.)  

 Also, $\sqrt{y^2-\frac{1}{16}} < \sqrt{y^2} = y$, so we have $x < y + z$, $y < z + x$ and $z < x + y$.  But these conditions are exactly those of the \href{https://artofproblemsolving.com/wiki/index.php?title=Triangle\_inequality}{triangle inequality}, so there does exist such a triangle.

 





\textbf{Solution 2 (Heron Bash)}

We can rewrite the equations as follows:

 
 \begin{align*} \frac{1}{4}x &= \sqrt{\left(y+\frac{1}{4}\right) \left(y-\frac{1}{4}\right)\left(\frac{1}{4}\right) \left(\frac{1}{4}\right)}+\sqrt{\left(z+\frac{1}{4}\right) \left(z-\frac{1}{4}\right) \left(\frac{1}{4}\right) \left(\frac{1}{4}\right)} \\ \frac{1}{5} y &= \sqrt{\left(z+\frac{1}{5}\right) \left(z-\frac{1}{5}\right) \left(\frac{1}{5}\right) \left(\frac{1}{5}\right)}+\sqrt{\left(x+\frac{1}{5}\right) \left(x-\frac{1}{5}\right) \left(\frac{1}{5}\right) \left(\frac{1}{5}\right)} \\ \frac{1}{6} z &= \sqrt{\left(x+\frac{1}{6}\right) \left(x-\frac{1}{6}\right) \left(\frac{1}{6}\right) \left(\frac{1}{6}\right)}+\sqrt{\left(y+\frac{1}{6}\right) \left(y-\frac{1}{6}\right) \left(\frac{1}{6}\right) \left(\frac{1}{6}\right)} \end{align*}


 Take the first equation. The first square root is the area of a triangle with side lengths $y,y,\frac{1}{2}$ by Heron’s Formula. Similarly, the second square root is the area of a triangle with side lengths $z,z,\frac{1}{2}$. If we connect these two triangles together at the $\frac{1}{2}$ side, we obtain a kite. The area of the kite is $\frac{1}{4}x$, and since the first diagonal is $\frac{1}{2}$, the second diagonal is $\frac{2\cdot\frac{1}{4}x}{\frac{1}{2}}=x$. If we draw this diagonal, we obtain two triangles with side lengths $x,y,z$. Let this triangle have area $A$. Then $\frac{1}{4}x=2A$; extend this for the other two equations to get $\frac{1}{5}y=2A$ and $\frac{1}{6}z=2A$. Therefore $x=8A$, $y=10A$, and $z=12A$. But we can find the area of the $x,y,z$ triangle with Heron's:
 \[A=\sqrt{(15A)(7A)(5A)(3A)}=15A^2\sqrt{7}\]Therefore $A=\frac{1}{15\sqrt{7}}$. Hence $x+y+z=30A=\frac{2}{\sqrt{7}}\Rightarrow\boxed{009}$.

 





\textbf{Solution 3 (Algebraic)}

Note that none of $x,y,z$ can be zero.

 Each of the equations is in the form 
 \[a=\sqrt{b^2-d^2}+\sqrt{c^2-d^2}\]

 Isolate a radical and square the equation to get 
 \[b^2-d^2=a^2-2a\sqrt{c^2-d^2}+c^2-d^2\] 

 Now cancel, and again isolate the radical, and square the equation to get 
 \[a^4+b^4+c^4+2a^2c^2-2a^2b^2-2b^2c^2=4a^2c^2-4a^2d^2\]

 Rearranging gives 
 \[a^4+b^4+c^4=2a^2b^2+2a^2c^2+2b^2c^2-4a^2d^2\]

 Now note that everything is cyclic but the last term (i.e. $-4a^2d^2$), which implies 
 \[-4x^2\cdot\frac1{16}=-4y^2\cdot\frac1{25}=-4z^2\cdot\frac1{36}\]

 Or

 
 \[x: y: z=4: 5: 6 \implies x=\frac{4y}5 \textrm{ and } z=\frac{6y}5\]

 Plug these values into the middle equation to get 
 \[\frac{256y^4+625y^4+1296y^4}{625}=\frac{800y^4}{625}+\frac{1800y^4}{625}+\frac{1152y^4}{625}-\frac{100y^2}{625}\]

 Simplifying gives 
 \[1575y^4=100y^2 \textrm{ but } y \neq 0 \implies y^2=\frac4{63} \textrm{ or } y=\frac2{3\sqrt7}\]

 Substituting the value of $y$ for $x$ and $z$ gives
 \[x+y+z = \frac{4y+5y+6y}5 = 3y = 3 \cdot \frac{2}{3\sqrt7} = \frac{2}{\sqrt7}\] 

 And thus the answer is $\boxed{009}$

 


 




\end{enumerate}